\documentclass[11pt,a4paper]{article}
\usepackage[T1]{fontenc}
\usepackage{lmodern,amsmath,amssymb,amsthm,mathtools,microtype}
\usepackage[margin=27mm]{geometry}
\usepackage[colorlinks=true,linkcolor=blue,citecolor=blue,urlcolor=blue]{hyperref}
\setlength{\emergencystretch}{2em}
\newtheorem{theorem}{Theorem}[section]
\newtheorem{proposition}[theorem]{Proposition}
\newtheorem{lemma}[theorem]{Lemma}
\newtheorem{corollary}[theorem]{Corollary}
\newcommand{\Q}{\mathbb Q}\newcommand{\Z}{\mathbb Z}
\DeclareMathOperator{\Res}{Res}\DeclareMathOperator{\Tr}{Tr}
\title{Curves and abelian families in an actual-trace square locus}
\author{}\date{17 November 2026}
\begin{document}\maketitle
\begin{abstract}
Fixing the rational input in the nine-contact derivative detector,
every algebraic curve in its guarded actual-trace square locus has
only finitely many rational points. There is no geometrically integral
curve defined over the rationals whose smooth projective normalization
has genus zero or one. The three elliptic point roles may vary
independently; no multiplication relation is needed. For a fixed root
point, a rational family parametrized by an abelian variety must also
have constant slope. These assertions restrict curve and abelian-family
constructions. They do not exclude individual rational witnesses,
constant-slope abelian families, or rational points varying on a
higher-dimensional nonabelian subvariety. The corresponding two-trace
general source has the same curve finiteness property; every rational
abelian family at a fixed input has constant common slope.
\end{abstract}

\section{The guarded locus at a fixed input}

Put
\[
 K=\Q(\sqrt2),\quad d=3-\sqrt2,\quad D=d^2,\quad
 E:y^2=x^3-Dx,\quad R=\Res_{K/\Q}E,
\]
and use the fixed point
\[
 P=\left(\frac{57-28\sqrt2}{8},
          \frac{-660+417\sqrt2}{32}\right).
\]
Fix $x_0\in\Q$ and an odd integer $B\ge3$. The three point roles
$A,H,U$ range independently over $R$. Write $f(Q)=1/x(Q)$,
$\zeta(Q)=-x(Q)/y(Q)$ and
\[
 r=\frac{\zeta(U)}{\zeta(H)},\qquad
 r'=\overline{\frac{\zeta(U-A)}{\zeta(H-P)}}.
\]
The six actual tests are
\[
 \begin{gathered}
 C_1=f(H),\quad C_2=f(U),\quad
 C_3=\overline{f(H-P)},\quad C_4=\overline{f(U-A)},\\
 C_5=r-r',\qquad C_6=r-x_0.
 \end{gathered}
\]
The bar interchanges the two geometric slots of the Weil restriction.
Set
\begin{equation}\label{eq:scale}
 \beta=x_0^2+B,\qquad T=\Tr x(A),\qquad
 M=\beta^4+T^4,\qquad a=M^4,
\end{equation}
and let $\mathcal S=\{-10,-8,-5,1,2,3,4,6,7\}$. The actual
detector functions are
\begin{equation}\label{eq:detector}
 \begin{gathered}
 F(a,z)=\prod_{c\in\mathcal S}((a+c)^2-z^2),\qquad
 u=-18F/F_a,\qquad w=u/(1+u)^2,\\
 D_i=C_i^2/(1+uC_i^2),\qquad
 W=\prod_{i=1}^6(1+D_i^2/u),\qquad
 V=\Tr(W^{-1}),\qquad \mathcal C=V^2+w(2-V).
 \end{gathered}
\end{equation}
Here $F_a$ is the partial derivative with $z$ fixed.

Let $\mathcal X_{B,x_0}$ be the locus $Y^2=\mathcal C$ on the
following open domain. Retain all point, affine-difference and quotient
guards of the derivative source \cite{Derivative,Criterion}. In
particular $A$ has finite abscissa in both slots; the tests $C_i$ are
finite; and $C_1,C_2,C_3,C_4$ are nonzero in both slots. Retain
the nonzero guards on $M$, $M^8-z^2$, $F$, $F_a$, $u$, $1+u$,
the balanced denominators, both slots of $W$, and $V$.
We may additionally remove the finitely many constant values of $a$
for which two numbers $\pm(a+c)$ coincide or one is zero.
Every rational point tuple satisfies
\begin{equation}\label{eq:large-a}
                    a\ge\beta^{16}\ge B^{16}>10,
\end{equation}
so this last removal discards no rational source point. There is no
nonzero condition on $C_5,C_6$ or the terminal ordinate $Y$.
Additional original guards or preparation inequalities only restrict
the locus further.

This is an algebraic locus in the actual point roles and the existing
slope. It has no independent trace parameter. The fixed-root locus
$A=[n]P$, $x_0=n\in\Z\setminus\{0\}$, and its sublocus
$U=[n]H$ are special cases. The results also apply after passing to
the source's finite point-chart covers.

\section{The fixed-point fibers}

\begin{lemma}\label{lem:fiber}
Fix a guarded rational point tuple $(A,H,U)$. The smooth projective
normalization of
\[
                         Y^2=\mathcal C(z)
\]
is geometrically integral of genus at least eight. It has eighteen
distinct rational branch points over $z=\pm(a+c)$ and two further
rational points over $z=\infty$, with $Y=\pm2$.
All twenty lie outside the original guarded affine source.
The guarded fiber has points over every completion of $\Q$.
\end{lemma}
\begin{proof}
At $z=\varepsilon(a+c)$, the distinctness supplied by $a>10$ gives
\[
 F=0,\qquad F_a\ne0,\qquad u=0,\qquad
                         \frac{du}{dz}=18\varepsilon.
\]
In each slot at least four tests are nonzero. A nonzero test contributes
\[
 \left(1+\frac{D_i^2}{u}\right)^{-1}
 =\frac{u(1+uC_i^2)^2}{u(1+uC_i^2)^2+C_i^4},
\]
which has a simple zero at $u=0$; a zero test contributes one.
Consequently $V=O(u^4)$ and
\[
             \mathcal C=2u+O(u^2),\qquad
             \frac{d\mathcal C}{dz}=36\varepsilon.
\]
These are eighteen simple rational zeros. The degree-two
function-field extension is geometrically nontrivial, and
Riemann--Hurwitz gives genus at least eight. This is the contact
calculation of \cite{Genus}; it requires no relation between $U$ and $H$.

Since $\sum_c c=0$, at infinity one has
\[
 F=-z^{18}+O(z^{16}),\qquad
 F_a=18az^{16}+O(z^{14}),\qquad u\sim z^2/a.
\]
Each nonzero test gives $D_i=u^{-1}+O(u^{-2})$ there, while a zero
test gives $D_i=0$. Thus $V=2+O(u^{-3})$ and
$\mathcal C=4+O(z^{-6})$. The local equation at infinity has
the two smooth rational points $Y=\pm2$, because its derivative in
$Y$ is nonzero at each. They are unramified. The eighteen contacts
violate $u\ne0$, and infinity is outside the affine slope chart.

At every finite place, taking $z$ sufficiently close to infinity makes
$\mathcal C$ a nonzero local square near four and keeps the finitely
many slope-dependent guards nonzero. At two, for example, require
$\mathcal C/4\in1+8\Z_2$. At the real place take $z>a+10$.
Then $u,w>0$, both slots of $W$ exceed one, and
$0<V<2$, so $\mathcal C>0$. These are points of the guarded
local fiber. This argument does not impose selected-contact prime
patterns or assert a global rational point in the guarded open set.
\end{proof}

\section{The rational simple abelian surface}

\begin{lemma}\label{lem:simple}
The surface $R$ is simple over $\Q$. Every positive power $R^s$
has no elliptic abelian subvariety defined over $\Q$.
\end{lemma}
\begin{proof}
The nonisogeny and simplicity statements are established in
\cite{Charts}; the following also verifies the field of definition.
Write $\sigma(d)=3+\sqrt2$ and
\[
 \delta=\frac{\sigma(d)}d=\frac{\sigma(d)^2}{7},\qquad
 \lambda=\frac{\sigma(d)}{\sqrt7},\qquad \lambda^2=\delta.
\]
The map $\phi(x,y)=(\lambda^2x,\lambda^3y)$ identifies
$E$ with $E^\sigma$ over $K(\sqrt7)$. The geometric
endomorphism ring of the characteristic-zero curve of invariant
$j=1728$ is $\Z[i]$, and all its endomorphisms are defined over
$K(i)$. The automorphism of $K(i,\sqrt7)$ fixing $K(i)$ and
negating $\sqrt7$ therefore negates every element of
\[
             \operatorname{Hom}_{\overline\Q}(E,E^\sigma)
                        =\phi\operatorname{End}_{\overline\Q}(E).
\]
It exists because $\sqrt7\notin K(i)$: the real subfield of
$K(i)$ is $K$, and $\sqrt7\notin K$. Hence
$\operatorname{Hom}_{K(i)}(E,E^\sigma)=0$.
The argument concerns rationality of isogenies, not geometric
nonisogeny. As a separate exact verification, at the degree-one place
of $K(i)$ above $17$ with $\sqrt2=6$, $i=4$, the two good
reductions have respectively $20$ and $16$ points.

If an elliptic abelian subvariety $J\subset R$ were defined over
$\Q$, its two projections after base change to
$R_K=E\times E^\sigma$ could not both be nonzero: two such
projections would be isogenies, contradicting the preceding
nonisogeny. Exactly one would be nonzero, making $J_K$ one of
the two factors. Galois interchanges those factors, so neither
descends to $\Q$. This proves simplicity. An elliptic subvariety
of $R^s$ would have a nonzero projection into some factor $R$,
again impossible.
\end{proof}

\begin{lemma}\label{lem:lowgenus}
Every morphism from a smooth projective geometrically integral
curve of genus at most one over $\Q$ to $R^s$ is constant.
No rational-point hypothesis on the curve is required.
\end{lemma}
\begin{proof}
A morphism $f:C\to R^s$ induces a homomorphism
$J(C)\to R^s$ through the difference map
$(c_1,c_2)\mapsto f(c_1)-f(c_2)$. This Jacobian construction
is defined over $\Q$ even if $C(\Q)$ is empty
\cite[Chapter III, Proposition 6.4]{Milne}. Its image is zero:
the Jacobian has dimension zero or one and
Lemma~\ref{lem:simple} excludes an elliptic image. Hence $f$ is
constant. Rational maps from the smooth projective curve to the
proper variety $R^s$ extend to morphisms, so the same conclusion
holds for them.
\end{proof}

\section{Finiteness on every fixed algebraic curve}

\begin{theorem}\label{thm:curves}
For fixed $B,x_0$ as above, no geometrically integral curve defined
over $\Q$, with smooth projective normalization of genus zero or
one, admits a nonconstant rational map into $\mathcal X_{B,x_0}$.
Every algebraic subset of $\mathcal X_{B,x_0}$ of dimension at
most one has only finitely many rational points. Both statements
hold for finite point-chart covers of this locus.
\end{theorem}
\begin{proof}
Suppose such a map from a curve of genus at most one exists.
Its projection to $R^3$ is constant by Lemma~\ref{lem:lowgenus}.
The constant point is rational, because the curve is geometrically
integral over $\Q$. Thus $(A,H,U)$ is fixed and $a>10$ by
\eqref{eq:large-a}. All four required reciprocal tests remain
finite and nonzero in both slots; vanishing of either final
comparison does not change Lemma~\ref{lem:fiber}.
If $z$ varies, the map induces a nonconstant map to the smooth
normalization of that fiber, whose genus is at least eight.
Riemann--Hurwitz rules this out. If $z$ is constant, $Y$ has
at most two values, so the original map is constant.

Every geometrically integral curve in the locus therefore has
normalization of genus at least two. Faltings's theorem makes
the rational points of that normalization finite \cite{Faltings}.
Normalization is an isomorphism away from finitely many singular
points, so these exceptions do not affect finiteness. If an
irreducible curve over $\Q$ is not geometrically integral, its
rational points lie in the finite intersections of distinct
conjugate geometric components. Zero-dimensional components also
contribute only finitely many points. Taking the finitely many
components proves the second assertion.

Under a finite point-chart cover, an algebraic curve cannot map
to a point. Its image is a curve of the preceding kind; the same
genus argument and finite fibers give both assertions.
\end{proof}

The theorem fixes the input. It allows $A,H,U$ to vary independently
and includes every torsion specialization that survives the displayed
guards. In particular it applies to the full fixed-root locus of
$\mathrm{TC}_B(n)$ and to its exact multiplication family. It does
not give a uniform bound as curves or inputs vary, or finiteness of
all rational points of the higher-dimensional locus.

\section{Nonconstant slopes in abelian families}

In this section fix the root point $A\in R(\Q)$ as well as
$B,x_0$, so that $a>10$ is a constant. Denote the resulting
guarded locus by $\mathcal X_{B,x_0,A}$. The numerator point $U$
is still independent of $H$.

\begin{theorem}\label{thm:abelian}
Let $J/\Q$ be an abelian variety. Every rational map
\[
                    J\dashrightarrow\mathcal X_{B,x_0,A}
\]
has constant slope $z$. In particular, a rational section after an
isogeny of the reference abelian surface cannot use a nonconstant
slope. No exclusion of constant-slope abelian families is asserted.
\end{theorem}
\begin{proof}
Suppose $z$ is nonconstant. The reference-point map to $R$ extends
to a morphism on $J$ and is a homomorphism followed by a translation
\cite[Chapter I, Theorem 3.2 and Corollary 3.7]{Milne}. By simplicity of $R$, it is
either constant or surjective. Work over $\overline\Q$ and let
$v$ be the order at any prime divisor of $J_{\overline\Q}$.

If the reference map is constant, its first and third tests are
nonzero constants, by the guards. If it is surjective, the common
zero locus of $C_{1,+},C_{3,+}$ is the inverse image of the point
\[
                            (O,\overline P)
                  \quad\text{in }R_{\overline\Q}=E\times E^\sigma.
\]
The common zero locus for the other slot is the inverse image of
$(P,O)$. Fibers of a surjective abelian homomorphism have constant
codimension two. Neither common zero locus contains a divisor.
We have used the exact fact that $f=1/x$ has its only zero at $O$.
Thus in either case and in each slot at least one of $C_1,C_3$
has valuation at most zero.

Suppose $\nu=v(u)>0$ is odd. For any nonzero test, parity gives
\[
 v(D_i)=\max(2v(C_i),-\nu),\qquad
 v(1+D_i^2/u)=\min(0,2v(D_i)-\nu)\le0.
\]
A zero test contributes the factor one. For a test of valuation at
most zero the last displayed valuation is at most $-\nu$.
Consequently, in both slots $v(W^{-1})\ge\nu$, and hence
$v(V)\ge\nu$. Since $v(w)=\nu$ and $2-V$ is a unit,
\[
                         v(\mathcal C)=\nu
\]
is odd. This contradicts the square identity $Y^2=\mathcal C$.
There is therefore no divisor on which $u$ has positive odd order.

At each of the eighteen constant contacts $z_c=\pm(a+c)$,
$u$ has a simple zero as a rational function of $z$. It follows
that every component of the zero divisor of $z-z_c$ has even
multiplicity. Therefore the rational function
\[
                            q=\prod_{z_c}(z-z_c)
\]
has even divisor on $J_{\overline\Q}$: the pole divisor is
eighteen times that of $z$, and its zero divisors are even.
The equation $s^2=q$ consequently defines an unramified
double cover of the smooth abelian variety. Explicitly, half of
$\operatorname{div}(q)$ defines a line bundle whose square is
trivial, and the associated square-root algebra is finite etale.
If $q$ is already a square, use the trivial cover. A connected
component $J'$ of this finite etale cover is an abelian variety
after choosing an origin. For a direct verification, embed the
ground field in $\mathbb C$ and write the complex abelian variety
as $\mathbb C^g/\Lambda$. A connected finite unramified cover
has the form $\mathbb C^g/\Lambda'$ for a finite-index sublattice
$\Lambda'$. It is projective, since the pullback of an ample
line bundle under the finite covering is ample. Hence this complex
torus is an abelian variety; this is all that the following
geometric contradiction requires.

The functions $z,s$ now give a nonconstant rational map from $J'$
to the fixed smooth hyperelliptic curve
\[
                          s^2=\prod_{z_c}(z-z_c),
\]
which has genus eight. This is impossible. Indeed compose with
its Abel--Jacobi embedding. The rational map into the Jacobian
extends to a morphism from $J'$ and is a homomorphism followed
by a translation. Its image, if one-dimensional, is an elliptic
translate and cannot equal a curve of genus eight. This contradiction
proves that $z$ is constant.
\end{proof}

Theorem~\ref{thm:abelian} uses the actual reciprocal trace, including
possible cancellations between its slots; such cancellation only
increases $v(V)$. It does not impose $U=[n]H$ or assume the last
two comparisons are nonzero. The fixed root and fixed input are
essential to the displayed constant-contact hyperelliptic argument.
The theorem leaves open a constant-slope abelian family, an isolated
rational witness, or a family on a nonabelian variety of dimension
at least two. Neither theorem disproves the eventual-positive
actual-trace realization criterion or supplies a new defining count.

\begin{corollary}\label{cor:classes}
The conclusions of Theorems~\ref{thm:curves} and \ref{thm:abelian}
also hold if $Y^2=\mathcal C$ is replaced by any one of
\[
 Y^2=(V+2)\mathcal C,\qquad
 Y^2=(V^2-V+2)\mathcal C,\qquad
 Y^2=(1+w(2-V))\mathcal C.
\]
They hold separately for every member of a fixed finite menu of
these three tests and the original test.
\end{corollary}
\begin{proof}
These are the three retained classes of \cite{Classes}. At every
contact their multipliers are respectively $2,2,1$, so each
radicand still has eighteen simple zeros. At a positive odd
divisorial valuation in the proof of Theorem~\ref{thm:abelian},
$v(V)\ge v(u)>0$ and $v(w)=v(u)>0$. Each multiplier is a
unit. The odd-order contradiction and the even-contact-fiber
argument are consequently unchanged. A rational map from an
irreducible variety into a finite union has image in one component,
so allowing a finite menu does not alter either conclusion.
\end{proof}

\section{The two-trace general source}

Use the general source of \cite[Section 1]{Reweighted} and its
two actual-trace equations from \cite{TwoTrace}. Fix its rational
input $X=x_0$ and odd coefficient exponent $L\ge3$. The scalar
$\rho=\rho_L(x_0)$ is now a fixed nonzero rational number.
There are nine independent point roles on $R^9$, and one common
slope $\lambda$. Keep all the displayed point, difference,
scalar, contact, weight and trace guards of that source.

Only the first detector is needed for the following assertion.
Its comparison scale is $a_0=1$. Put
\[
 \begin{gathered}
 \Theta=\{0,-1/2,-1,1/2399,-32/407\},\qquad
 \Theta_0=(2+\Theta)\cup(4+\Theta),\\
 \delta_0=\lambda+1,\qquad
 R_0=\frac{\prod_{\theta\in\Theta_0}(\lambda-\theta)}
             {\delta_0^{10}},\qquad
 u_0=\rho R_0,\qquad w_0=\frac{u_0}{(1+u_0)^2}.
 \end{gathered}
\]
Its ten distinct rational contacts are disjoint from the pole $-1$.
With the field tests and scalar tests of \cite{Reweighted},
\[
 \begin{split}
 G_{0,k}&=R_0^{\kappa_k+1}\frac{R_0S_k}{1+R_0S_k},\\
 W_0&=\prod_{i=1}^{16}
       \left(1+\frac{C_i^2}{\delta_0^2w_0^{2a_i-1}}\right)
       \prod_{k=1}^{6}
       \left(1+\frac{G_{0,k}^2}{w_0^{2\kappa_k-1}}\right),\\
 V_0&=\Tr(W_0^{-1}),\qquad
 \mathcal C_0=V_0^2+w_0(2-V_0).
 \end{split}
\]
The orders $a_i$ are one except for $i=12,13,14,16$, where
they are three; $\kappa_k=(1,1,1,1,3,3)$.
The first ten tests are
$(f(U_i))_{i=0}^{4}$ and
$(\overline{f(U_i-Q_i)})_{i=0}^{4}$, with $U_0=H$, $Q_0=P$.
Only the nonzero reference tests $f(H)$ and
$\overline{f(H-P)}$ are needed below. All other finite tests may
vanish, including any zero reciprocal numerator admitted by a native
point chart. The second detector and its actual square equation
$Y_1^2=\mathcal C_1$ are retained in full.

\begin{theorem}\label{thm:general}
In this fixed-input two-trace locus, every algebraic subset of
dimension at most one has finitely many rational points, and no
geometrically integral curve over $\Q$ of normalization genus
zero or one admits a nonconstant rational map into the locus.
Every rational map from an abelian variety over $\Q$ into the
locus has constant common slope $\lambda$. These conclusions
also hold for its finite point-chart covers and separately for
the small-class alternatives of Corollary~\ref{cor:classes}.
\end{theorem}
\begin{proof}
First fix all nine guarded point roles. At each contact
$\lambda=\theta$, $w_0$ has a simple zero, $\delta_0$ is
nonzero. In each slot the nonzero reference reciprocal supplies a
field factor with a simple pole. All other field factors either
have a pole or equal one. Every scalar test is a finite constant, so
$G_{0,k}=O(R_0^{\kappa_k+2})$ and its factor equals
$1+O(R_0^5)$. It follows that $V_0=O(w_0)$ and
\[
                           \mathcal C_0=2w_0+O(w_0^2).
\]
The normalization of $Y_0^2=\mathcal C_0(\lambda)$ therefore
has ten simple rational branch points and genus at least four.
Any curve component in the two-square fiber maps finitely to
this first cover, so its normalization has genus at least four.
The projection to $R^9$, Lemma~\ref{lem:lowgenus}, and Faltings's
theorem now prove the first two assertions exactly as in
Theorem~\ref{thm:curves}, including singular and non-geometric
components.

For an abelian family, consider a divisor where
$\lambda-\theta$ has positive odd order $\nu$. The scalar
$\rho$ and $\delta_0$ are units there, and $v(w_0)=\nu$.
Every detector factor has nonpositive valuation: its nonconstant
summand has valuation an even integer minus an odd multiple of
$\nu$, and therefore cannot have valuation zero. Zero tests
again contribute one. In each slot at least one of $f(H)$ and
$\overline{f(H-P)}$ has valuation at most zero, by the same
constant-or-surjective reference-map argument used in
Theorem~\ref{thm:abelian}. Its field factor has valuation at
most $-\nu$. Hence $v(V_0)\ge\nu$ and
$v(\mathcal C_0)=\nu$, contradicting rational squareness.

Every zero divisor of $\lambda-\theta$ is consequently even.
The product of these ten functions has even divisor, since its
pole divisor is ten times that of $\lambda$. Its square root on
a finite etale abelian cover would give a nonconstant map to the
fixed genus-four curve
\[
                           s^2=\prod_{\theta\in\Theta_0}
                                            (\lambda-\theta).
\]
The same Jacobian argument forces $\lambda$ to be constant.
The small-class multipliers are units at the contact valuations,
so their inclusion does not change the proof.
\end{proof}

The first comparison scale is constant, so the abelian-family
assertion in Theorem~\ref{thm:general} does not require fixing
any root point. It concerns the displayed reweighted general
source and its actual two-square replacement. It asserts neither
failure of its realization criterion nor a lower bound for other
Diophantine constructions.

\section{Exact algebra checks}
The companion checker verifies the conjugate-twist isomorphism,
the two point counts at $17$, the eighteen contact derivatives,
and the expansion at infinity. It also checks the two universal
valuation identities used in Theorem~\ref{thm:abelian} on a
finite set of integer boundary cases. The general valuation,
Jacobian, and etale-cover arguments are proved above; they are
not inferred from those finite checks.

\begin{thebibliography}{9}\raggedright
\bibitem{Derivative}\emph{An eleven-variable construction under rational
$abc$ using a derivative detector},
\nolinkurl{papers/conditional_few_variables/derivative_detector11_20261117_v63.tex}.
\bibitem{Criterion}\emph{One actual-trace square in a derivative detector},
\nolinkurl{papers/pole_parameters/single_trace_derivative_criterion_20261117_v64.tex}.
\bibitem{Genus}\emph{Genus of the actual-trace fibers in a derivative detector},
\nolinkurl{negative_results/papers/pole_parameters/derivative_trace_genus_20261117_v65.tex}.
\bibitem{Charts}\emph{An obstruction to one-parameter presentations of an elliptic point},
\nolinkurl{negative_results/papers/elliptic_charts/elliptic_chart_obstruction_20261107_v7.tex}.
\bibitem{Classes}\emph{Small polynomial classes for the actual detector trace},
\nolinkurl{papers/pole_parameters/small_trace_classes_20261117_v58.tex}.
\bibitem{Reweighted}\emph{Quadratic reweightings of native elliptic detectors},
\nolinkurl{papers/pole_parameters/reweighted_elliptic_detectors_20261116_v53.tex}.
\bibitem{TwoTrace}\emph{A rational conic criterion using the actual detector trace},
\nolinkurl{papers/pole_parameters/trace_conic_criterion_20261117_v57.tex}.
\bibitem{Milne} J. S. Milne, \emph{Abelian Varieties}, version 2.0,
2008, Chapters I and III.\quad
\url{https://www.jmilne.org/math/CourseNotes/AV.pdf}.
\bibitem{Faltings} G. Faltings,
\emph{Endlichkeitss\"atze f\"ur abelsche Variet\"aten \"uber Zahlk\"orpern},
Invent. Math. \textbf{73} (1983), 349--366; erratum,
\textbf{75} (1984), 381.\quad
\url{https://doi.org/10.1007/BF01388432}.
\end{thebibliography}
\end{document}
